\documentclass[12pt]{article}

\usepackage{geometry}
\usepackage{titling}
\usepackage{amsmath,amscd}
\usepackage{fontspec}
\usepackage{unicode-math}
\usepackage{pifont}
\usepackage{xcolor}
\usepackage{colortbl}
\usepackage[dvipsnames]{xcolor}
\usepackage{listings}
\usepackage{minted}
\usepackage{tikz}
\usepackage{nicematrix}
\usepackage[most]{tcolorbox}

\usetikzlibrary{fit}

\geometry{a4paper, top=2.5cm, bottom=2.5cm, left=2.6cm, right=2.5cm}
\setlength{\droptitle}{-3cm}

\setmonofont{DejaVu Sans Mono}
\setmainfont{Libertinus Serif}
\setsansfont{Libertinus Sans}
\setmathfont{STIX Two Math}

\definecolor{lightyellow}{rgb}{1, 1, 0.8} % Define a subtle highlight color
\definecolor{codepy}{rgb}{0.9, 0.85, 0.85} %{#F8BC65}
\definecolor{codegray}{rgb}{0.5,0.5,0.5}
\definecolor{codecoffee}{rgb}{0.75,0.6,0.3}
\definecolor{codecomment}{rgb}{0.97, 0.74, 0.4}
\definecolor{lightblue}{rgb}{0.55,0.75,0.9}
\definecolor{midgray}{rgb}{0.4, 0.4, 0.3}
\definecolor{indexwhite}{rgb}{0.9, 0.9, 0.9}
\definecolor{darkgold}{HTML}{8F3F0F}
\definecolor{amber}{RGB}{255, 191, 0}
\definecolor{brickred}{RGB}{205, 92, 92}

%\newfontfamily\symbolfont{DejaVu Sans}

\newcommand{\srccodesize}{\fontsize{9pt}{10pt}\selectfont}

%\newcommand{\warnicon}{
%    {\symbolfont\textcolor{red}{\textbf\Large{\symbol{"26A0}}}}%
%}


\lstdefinestyle{mypystyle}{
    aboveskip=2pt,
    belowskip=2pt,
    frame=single,
    framerule=0pt, % disable the frame rules
    framextopmargin=2pt,
    framexbottommargin=2pt,
    lineskip=3pt,
    % expand the margins to pull everything inside the box
    framexleftmargin=4.4mm,     % Makes the left side of the black box wider
    xleftmargin=5.8mm,
    xrightmargin=4pt,
    commentstyle=\color{codecomment},
    keywordstyle=\color{lightblue},
    numberstyle=\srccodesize\color{white!40},
    stringstyle=\color{codecoffee},
    basicstyle=\ttfamily\srccodesize\color{codepy},
    breakatwhitespace=false,
    breaklines=true,
    % captionpos=b,
    keepspaces=true,
    numbers=left,
    numbersep=0pt,
    firstnumber=1,  % Forces reset to 1 for EVERY listing
    showspaces=false,
    showstringspaces=false,
    showtabs=false,
    tabsize=2,
    language=Python,
    backgroundcolor=\color{black!80},
    emph=[2]{__init__, calculate},
    emphstyle=[2]\color{funcColor},
}
\lstset{style=mypystyle}

\newcommand{\vsp}{\vspace*{8pt}}
\newcommand{\snl}{\vspace*{2pt}}

\title{Introduction to NumPy Arrays and Matrices}
\author{Devi Prasad M}
\begin{document}

% A soft, light gray background that is easy on the eyes
\pagecolor{yellow!30!brown!20}

% A rich, warm cream that mimics expensive book paper
% \pagecolor{yellow!20!brown!10!white!70}

% A sophisticated, cool-toned paper look
% \pagecolor{cyan!10!yellow!10!black!6}

% A soft, warm sepia that feels gentle on the eyes
% \pagecolor{red!10!yellow!18!black!6}

% High-contrast but low-glare dark mode text theme
%\pagecolor{yellow!5!cyan!10!black!85}
%\color{white!90!yellow!10} % Pair with this text color

\maketitle

\section{Working with ndarray}
An \emph{ndarray} object is the core data structure in NumPy, and
it can represent vectors, matrices, and higher-dimensional arrays.
NumPy arrays are \emph{fixed-size}, homogeneous data structures.

\vsp\subsection{Creating a 1D array}
\begin{lstlisting}[language=Python]
    a = np.array(np.arange(1, 25, 1, np.int8))
    assert(a.shape == (24,))
    assert(a.size == 24)
    assert(len(a) == 24)
    assert(a.ndim == 1)
\end{lstlisting}

$ndim$ represents the number of dimensions (axes) of
the array. In this case, since $a$ is a 1D array.

The $len$ function returns the number of elements along the
first axis of an array. For a 1D array, this is simply the total
number of elements in the array.

\subsection{Reshaping a 1D array into a 2D array}
The \emph{reshape} method is used to change the shape of
an array without changing its data. The total number of elements
in the reshaped array must match the total number of elements
in the original array.

\vsp\begin{lstlisting}[language=Python, escapeinside={(*@}{@*)}]
    a = np.array(np.arange(1, 25, 1, np.int8))
    assert (a == [x for x in range(1, 25)]).all()
    b = (*@\setlength{\fboxsep}{0pt}\colorbox{darkgold!90}{a.reshape}@*)(2, 12)
    assert(b.shape == (2, 12))
    assert(b.size == 24)
    (*@\setlength{\fboxsep}{0pt}\colorbox{darkgold}{assert(len(b) == 2)}@*)
    assert(b.ndim == 2)
\end{lstlisting}

In the following we create a 2D array with 4 rows and 6 columns.
The total number of elements in the reshaped array must match
the total number of elements in the original 1D array.

\vsp\begin{lstlisting}[language=Python, escapeinside={(*@}{@*)}]
    a = np.array(np.arange(1, 25, 1, np.int8))
    assert (a == [x for x in range(1, 25)]).all()
    c = a.reshape(4, 6)
    assert(c.shape == (4, 6))
    assert(c.size == 24)
    (*@\setlength{\fboxsep}{0pt}\colorbox{darkgold}{assert(len(c) == 4)}@*)
    assert(c.ndim == 2)
\end{lstlisting}

\subsection{Accessing rows and columns in a 2D array}
\vsp\begin{lstlisting}[language=Python, escapeinside={(*@}{@*)}]
    a = np.array(np.arange(1, 25, 1, np.int8))
    b = a.reshape(2, 12)
    # Check the first row of the reshaped array
    assert (b[0] == [x for x in range(1, 13)]).all()
    # Check the second row of the reshaped array
    assert (b[1] == [x for x in range(13, 25)]).all()
    # Check the first column of the reshaped array
    assert (b[:, 0] == [1, 13]).all()
    assert np.array_equal(b[:, 0], [1, 13])
    # access the second column
    assert np.array_equal(b[:, 1], [2, 14])
    # the last column
    assert np.array_equal(b[:, 11], [12, 24])
\end{lstlisting}

\subsection{Reshaping the 2x12 array into a 2x3 array}
We can further reshape the 2x12 array into a 2x3 array by slicing the first
three columns of the 2x12 array.

\vsp\begin{lstlisting}[language=Python, escapeinside={(*@}{@*)}]
    a2x3 = b[:, 0:3]
    assert(a2x3.shape == (2, 3))
    assert(a2x3.size == 6)
    (*@\setlength{\fboxsep}{0pt}\colorbox{darkgold}{assert(len(a2x3) == 2)}@*)
    assert(a2x3.ndim == 2)
    # access the last column of the 2x3 array
    assert np.array_equal(a2x3[:, 2], [3, 15])
\end{lstlisting}

\section{The Memory of \emph{ndarray} Objects}
NumPy arrays enable flexible reshaping and slicing of the underlying data
without actually copying data. When an array is sliced or reshaped,
NumPy creates a \emph{view} of the same \emph{base} array.
When a new view is created, it does not copy the data but instead provides
a new way to access the underlying array. This strategy avoids unnecessary
data duplication for large arrays.

In the following exmaple, notice how modifying the view affects the original array.
\vsp\begin{lstlisting}[language=Python, escapeinside={(*@}{@*)}]
    a = np.array(np.arange(1, 25, 1, np.int8))
    a_slice = a[10:20] # holds a[10]...a[19]
    assert a_slice.size == 10 and a_slice.shape == (10,) and a_slice.ndim == 1

    ### create a view of the original array
    assert a.base is None  # 'a' is the base array
    assert a_slice.base is a  # 'a_slice' is a view of 'a'
    # Modifying 'a_slice' will modify 'a'
    a_slice[0] = 127 # modifies a[10]
    assert a[10] == 127
    assert a_slice.base is a
\end{lstlisting}

We can also create a view of a view. In this case, the new view will still reference
the original base array, not the intermediate view. This means that changes made
to the new view will affect the original array, not the intermediate view.

\vsp\begin{lstlisting}[language=Python, escapeinside={(*@}{@*)}]
    ### mint a view of a view
    c = a_slice.reshape(2,5)
    assert c.base is a  # 'c' is a view of 'a', NOT 'a_slice'!
    c[0, 0] = 127
    assert a_slice[0] == 127
    assert a[10] == 127
\end{lstlisting}


\section{Combining NumPy arrays}
Two \emph{ndarray} objects can be combined using the \emph{concatenate} function.
This function takes a sequence of arrays and joins them along an existing axis.
The arrays must have the same shape, except in the dimension corresponding
to the axis along which they are concatenated.

\vsp\begin{lstlisting}[language=Python, escapeinside={(*@}{@*)}]
    a = np.concatenate(([1,2,3], [4,5], [7], [8,9], []), axis=0)
    assert(np.array_equal(a, [1, 2, 3, 4, 5, 7, 8, 9]))
\end{lstlisting}

\begin{tcolorbox}[
    colframe=brickred,
    colback=white,
    title=Efficiency Note,
    fonttitle=\bfseries
]
NumPy arrays have a fixed size: once an array is created, its shape cannot be
changed in place. Therefore, operations such as repeated
concatenation are inefficient - each concatenation creates a new array and
copies the accumulated data. When building an array incrementally, it is
more efficient to first collect the data in a Python list and create a
NumPy array from that list.
\end{tcolorbox}

\end{document}